%!TEX root = main.tex
%=========================================================

\section{Related Work}
\label{sec:related}

\y{this section lacks a cohesive narrative. it should identify a gap, and relationships to similar works}
% We review related work on broadcast in MANETs, adaptation mechanisms in wireless and mobile networks, and studies of mobility in these systems.

A survey by Ruiz and Bouvry~\citep{ruiz2015survey} provides a comprehensive taxonomy of controlled flooding and overlay-based protocols, also covered in surveys by Reina \emph{et~al.}~\citep{Reina2015} and Santi~\citep{Santi:2005:TCW:1089733.1089736}.
%
% Broadcast protocols consider either fixed- or a variable-range communication~\citep{woon2011variable,park2013efficient}.
% Our focus in this paper is on fixed-range communication.
Our previous work~\citep{DensityMobDrivenEval} and the work of Williams and Camp~\citep{Williams2002} present comparative evaluations of broadcast protocols from different classes.
Hess \emph{et~al.}~\citep{Hess:2015:DHM:2856149.2840722} survey human mobility models for mobile networking.
Levy walk are a \emph{general} model~\citep{rhee2011levy}.
Other proposals allow modeling mobility in specific scenarios, such as disaster areas~\citep{Aschenbruck2009}.

The \emph{broadcast storm} problem associated with flooding is identified by Ni et al.~\citep{BroadcastStormProblem}, who also propose \emph{controlled flooding} using counters and timers as we use in this paper.
The collection of topology information~\citep{ProbFlood} or GPS locations~\citep{abba} allows the optimizing of controlled flooding. However, they are costlier and more sensitive to high mobility.

Overlay-based broadcast is based on the construction of connected dominating sets (CDS)~\citep{stojmenovic-cds}, with the goal of minimizing the number of necessary relay nodes.
MPR~\citep{mpr} or TSS~\citep{Nikolov2015} create CDS by guaranteeing the coverage of two-hop neighbors by at least one relay node.
Ahn and Lee~\citep{mpr-improvement} improve MPR to guarantee that each 2-hop neighbor is covered by $m$ such relays.
Other approaches~\citep{woon2011variable,park2013efficient} consider variable-range local broadcast to optimize the CDS.

Determining the optimal configuration for a MANET broadcast algorithm is difficult.
% Adaptation techniques have been proposed for different protocols and application contexts.
Kokuti and Vilmos~\citep{kokuti2012adaptive} propose adaptation techniques to dynamically decide on parameters in controlled flooding.
Tseng, Ni and Shih~\citep{2001:AAR:876878.879323} similarly propose the ACF adaptation mechanism that we used in our evaluation (\S\ref{sec:evaluation}).
Controlled flooding adaptation can use measures of node velocity, as for S-H flooding~\citep{adaptive_group_com_iscc02} detailed in~\S\ref{sec:evaluation}, or measures of observed contention and collisions~\citep{2001:AAR:876878.879323}.
Liarokapis and Shahrabi~\citep{prob-based-adaptive-tech} propose techniques to switch from simple to controlled flooding, and adapt switching threshold to locally measured density.
Artimy~\citep{artimy2007local} proposes density estimation and dynamic range adaptation for vehicular networks, distinguishing between free-flow and congested phases.

The term \emph{emergent overlay} has been previously used in another context, wired unstructured P2P network~\citep{galuba2007generic}, to denote the \emph{on demand} construction of distributed hash tables~\citep{jelasity2009t} following a centralized decision.
Bellavista \emph{et al.}~\citep{Bellavista2013} combine a static sensor network using an overlay and a MANET with controlled flooding for data collection for an IoT scenario, but the use of either protocol is determined in advance.
We are not aware of any previous approach combining in a single MANET controlled flooding and overlays.


% broadcast algorithms in manets
% \begin{itemize}
% 	\item Surveys:
% 	\begin{itemize}
% 		\item : 2015 survey with the taxonomy
% 		\item \citep{Reina2015}: 2015 survey on probabilistic broadcast approaches
% 		\item \citep{Santi:2005:TCW:1089733.1089736}: 2005 survey on topology control in manets
% 	\end{itemize}
% 	\item Controlled flooding:
% 	\begin{itemize}
% 		\item \citep{BroadcastStormProblem}: 3 controlled flooding approaches to deal with broadcast storms
%
% 		% \item \citep{ProbFlood}: 2003 broadcast approach to privilege that nodes located at the border of transmissions range act as relays
% 	\end{itemize}
% 	\item Overlay-based
% 	\begin{itemize}
% 		\item \citep{stojmenovic-cds} 1st paper in determining connected dominating sets (CDS)
% 		\item \citep{Nikolov2015} 2015 a more recent approach to build CDS
% 		\item \citep{mpr} another approach to create CDS
% 		\item \citep{mpr-improvement}: 2014 recent improvement of MPR (multi-point relaying technique) to build backbones
% 		\item \citep{overlays_for_manets} 2015 recent study about how to maintain P2P overlays in MANETs for routing
% 	\end{itemize}
%
% 	%		\item \citep{clustering-approach}: 2009 approach for clusters within MANETs that uses location and energy metrics \er{for clusters? from the title it seems it uses broadcast as a tool, so you could rather use it in the intro to motivate the uses of broadcast}
%
% \end{itemize}

% studies of broadcast performance
% \begin{itemize}
% 	\item \citep{DensityMobDrivenEval} our evaluation study of broadcast over different scenarios of mobility and density
% 	\item evaluation of broadcast protocols~\citep{Williams2002}, there is also a first classification of protocols
% \end{itemize}

% adaptation mechanisms in wireless systems, manets, sensor networks
% \begin{itemize}
% 	\item \citep{artimy2007local} seems to also use overlays and flooding in vehicular networks
% 	\item \citep{kokuti2012adaptive}
% 	\item \citep{adaptive_group_com_iscc02} is basically one adaptation policy where the switching criteria is based on the relative velocity of nodes ; nodes switch between a list of controlled flooding approaches
% 	\item in this work \citep{Bellavista2013} overlays are build on the top of wireless sensor networks (WSN) to speed up data dissemination. Such data is classified with a certain priority; interoperability issues between WSN and MANETs are tackled but no between dissemination protocols
% 	\item 2003 \citep{2001:AAR:876878.879323} this work (aprox 700 citaitons) follows the paper about the broadcast storm problem to point out the need of adapting the value of thresholds shown in \citep{BroadcastStormProblem}
% 	\item \citep{hypergossiping} : adaptive approach based on gossiping
% 	\item \citep{prob-based-adaptive-tech} thresholds are changed on the fly without using control messages (seems promising, still have to read carefully)
% 	\item \citep{whitbeck2010hymad}: adaptation for routing in dense and dynamic networks
% \end{itemize}

% mobility traces and models
% \begin{itemize}
% 	\item \citep{Hess:2015:DHM:2856149.2840722}: a recent survey
% 	\item \citep{Aschenbruck2009}: mobility model for disaster area scenarios
% 	\item \citep{rhee2011levy}: Levy walk nature, classical paper
% \end{itemize}
%
% emergent overlays elsewhere
% \begin{itemize}
% 	\item This paper~\citep{galuba2007generic} uses the term in another context, an unstructured p2p network where a DHT or another structure is bootstrapped on demand
% \end{itemize}
